\section*{第 \ref{ch:g8:the-atmosphere} 章——大气}

\begin{solution}{exo:g8:the-atmosphere:1}
(a) \omterm{def:g8:the-atmosphere:layers}{平流层}；(b) \omterm{def:g8:the-atmosphere:layers}{对流层}；(c) \omterm{def:g8:the-atmosphere:layers}{中间层}；(d) \omterm{def:g8:the-atmosphere:layers}{热层}。
\end{solution}

\begin{solution}{exo:g8:the-atmosphere:2}
(a) \qty{15}{\degreeCelsius}。(b) 约 \qty{-56}{\degreeCelsius}。
(c) 约 \qty{-3}{\degreeCelsius}。(d) 约 \qty{-86}{\degreeCelsius}。
(e) 从 \qty{11}{km} 到约 \qty{20}{km}（以及在
\qty{50}{km} 附近和刚高于 \qty{86}{km} 处的几千米范围内）。
\end{solution}

\begin{solution}{exo:g8:the-atmosphere:3}
(a) 错误：温度在\omterm{def:g8:the-atmosphere:layers}{对流层}中下降，在\omterm{def:g8:the-atmosphere:layers}{平流层}中上升，在\omterm{def:g8:the-atmosphere:layers}{中间层}中下降，
并在\omterm{def:g8:the-atmosphere:layers}{热层}中再次上升。(b) 正确。
(c) 错误：它位于\omterm{def:g8:the-atmosphere:layers}{平流层}。(d) 正确。
\end{solution}

\begin{solution}{exo:g8:the-atmosphere:4}
(a) \omterm{def:g8:the-atmosphere:layers}{对流层}。(b) \omterm{def:g8:the-atmosphere:layers}{平流层}（标准大气中
\omterm{def:g8:the-atmosphere:layers}{对流层顶}位于 \qty{11}{km} 处）。(c) \omterm{def:g8:the-atmosphere:layers}{平流层}。(d) \omterm{def:g8:the-atmosphere:layers}{中间层}。
\end{solution}

\begin{solution}{exo:g8:the-atmosphere:5}
(a) $20 - 6.5 \times 2 = 20 - 13 = \qty{7}{\degreeCelsius}$。(b) $20 -
6.5 \times 3 = 20 - 19.5 = \qty{0.5}{\degreeCelsius}$。(c) 空气需要
降温 20 度，$20 \div 6.5 \approx 3.1$，故约在
\qty{3.1}{km} 处。
\end{solution}

\begin{solution}{exo:g8:the-atmosphere:6}
(a) $3 - 0.5 = \qty{2.5}{km}$。(b) $6.5 \times 2.5 = 16.25$，约降低 16
度。(c) $18 - 16.25 \approx \qty{2}{\degreeCelsius}$。
\end{solution}

\begin{solution}{exo:g8:the-atmosphere:7}
(a) \qty{500}{hPa}、\qty{250}{hPa} 和 \qty{125}{hPa}。(b) 靠近
地面时，该规则较为准确（\qty{500}{hPa} 对 \qty{506}{hPa}）。在更高处
该规则给出的结果偏大：\qty{250}{hPa} 对 \qty{227}{hPa}，以及
\qty{125}{hPa} 对 \qty{96}{hPa}。减半规则只是一个
粗略的估计。
\end{solution}

\begin{solution}{exo:g8:the-atmosphere:8}
(a) $10~\text{hPa} = 10^1$~hPa；$10^3 \div 10^1 = 10^2$：小一百
倍。(b) $0.01~\text{hPa} = 10^{-2}$~hPa。(c) $10^3 \div 10^{-2} =
10^5$：小十万倍。
\end{solution}

\begin{solution}{exo:g8:the-atmosphere:9}
(a) $6 \div 1.5 = 4$，宽 4 倍。(b) 气球上升时，周围空气的
压强越来越小，因此气球内的气体膨胀，气球
不断变大，直至球皮破裂。(c) 气球破裂后，降落伞
可减缓探空仪的下落速度，使其不会伤人或
损坏任何物品。
\end{solution}

\begin{solution}{exo:g8:the-atmosphere:10}
(a) $20 \div 7 \approx 2.9$，厚约 2.9 倍。(b) 靠近两极时：
\qty{11}{km} 比 \qty{7}{km} 高 \qty{4}{km}，却比 \qty{20}{km} 低 \qty{9}{km}（该数值描述的是中\omterm{def:g6:latitude-longitude:latitude}{纬度}，而非\omterm{def:g3:continents-and-oceans:equator}{赤道}地区）。
(c) 开始时，\omterm{def:g8:the-atmosphere:layers}{对流层顶}约在 \qty{7}{km} 处，
低于飞机所在高度，飞机位于\omterm{def:g8:the-atmosphere:layers}{平流层}。最后，\omterm{def:g8:the-atmosphere:layers}{对流层顶}约在
\qty{20}{km} 处，飞机位于\omterm{def:g8:the-atmosphere:layers}{对流层}。
\end{solution}

\begin{solution}{exo:g8:the-atmosphere:11}
(a) $1 - 506 \div 1013 \approx 0.50$：\textbf{50~\%}。(b) $1 - 227 \div
1013 \approx 0.78$：\textbf{78~\%}。(c) $1 - 12 \div 1013 \approx 0.99$：
\textbf{99~\%}。(d) 约 1~\%。
\end{solution}

\begin{solution}{exo:g8:the-atmosphere:12}
(a) $6371 \div 15 \approx 425$~千米，即每厘米相当于 425 千米。(b) $100 \div 425
\approx 0.24$~厘米，约 \qty{2.4}{mm}。(c) $11 \div 425 \approx
0.026$~厘米，约 \qty{0.26}{mm}。(d) $0.26 \div 0.1 \approx 2.6$：约
两三张纸。
\end{solution}

\begin{solution}{pb:g8:the-atmosphere:1}
\textbf{1.} $300 \times 10 = 3000$~米。

\textbf{2.} $6000 \div 300 = 20$ 分钟；$\num{30000} \div 300 = 100$
分钟。

\textbf{3.} 100 分钟，即 1 小时 40 分钟。

\textbf{4.} 在 \qty{8}{km} 处：\omterm{def:g8:the-atmosphere:layers}{对流层}（气温仍在
下降）。在 \qty{30}{km} 处：\omterm{def:g8:the-atmosphere:layers}{平流层}。

\textbf{5.} $15 - (-50) = 65$ 度。

\textbf{6.} $65 \div 10 = \qty{6.5}{\degreeCelsius}$ 每千米。

\textbf{7.} 从 \qty{12}{km} 到 \qty{20}{km}，气温保持在
\qty{-56}{\degreeCelsius}。降温在 \qty{10}{km} 与
\qty{12}{km} 之间停止：\omterm{def:g8:the-atmosphere:layers}{对流层顶}约位于 \qty{11}{km} 处。

\textbf{8.} 气温上升：\qty{22}{km} 处为 \qty{-55}{\degreeCelsius}，
\qty{30}{km} 处为 \qty{-47}{\degreeCelsius}。气球位于
\omterm{def:g8:the-atmosphere:layers}{平流层}，被吸收太阳紫外线的臭氧
加热。

\textbf{9.} $1013 \div 2 \approx \qty{507}{hPa}$；该气压出现在
\qty{4}{km}（\qty{617}{hPa}）与 \qty{6}{km}（\qty{472}{hPa}）之间。

\textbf{10.} $1013 \div 4 \approx \qty{253}{hPa}$。

\textbf{11.} 介于 \qty{10}{km}（\qty{265}{hPa}）与
\qty{12}{km}（\qty{194}{hPa}）之间。

\textbf{12.} 从 \qty{10}{km} 处起，气压还需再下降 $265 - 253 =
12$~hPa，该段在 2~千米内共下降 71~hPa：$12 \div 71 \times 2
\approx 0.34$~千米。此时气压为\omterm{def:g4:world-physical-map:altitude}{海平面}值的四分之一，所在高度约
\textbf{\qty{10.3}{km}}。
\end{solution}
